\scanpage{135}
\mspage{67}

We shall therefore assume that (37) does indeed define a structure of extensions, and we shall have a homomorphism of exact sequences

\diagram{\begin{tikzcd}[ampersand replacement=\&,row sep=large,column sep=large]
\deleted{(38)} \& 1 \arrow[r,dash] \& T^J \arrow[r] \arrow[d] \& \mathcal E_{U_{J,\alpha}}^{\bar L} \arrow[r] \arrow[d] \& \mathcal E_{D_{J,\alpha}}^{\tilde D_{J,\alpha}} \arrow[r] \arrow[d] \& 1 \\
(38) \& 1 \arrow[r,dash] \& \pi_{U/K}^{\bar L} \arrow[r] \& \mathcal E_U^{\bar L} \arrow[r] \& {\unsure{\mathcal E|\bar{\mathcal E}}}_K^{\bar K} \arrow[r] \& 1
\end{tikzcd}}

--- where we have \unsure{corrected|modified} our aim, noting that (37) only makes sense as it stands if $D_J$ is connected, \inserted{i.e. $X_J$ connected} that otherwise one must look at the connected components $X_{J,\alpha}$ of $X_J$, inducing the connected components $D_{J,\alpha}$ of $D_J$, and $U_{J,\alpha}$ of $U_J$, and write the exact sequence for each component. Of course, there are more intrinsic things: write with the \deleted{étale \unsure{groupoid|Galois|geometric|profinite}}\inserted{multi-Galois} topoi, and their fundamental groupoids....\deleted{\unsure{classes of} $\mathcal E_U^{\bar L}$ \unsure{conjugacy|conjugates}}

We thus see that the homomorphisms
\[
\tag{39}
T^J\longrightarrow\pi_{U/K}^{\bar L}
\]
induced by the different $x\in D_J$ --- each being in fact, when one returns to the precise definition, only defined modulo \unsure{opérat|operations} \unsure{direct|direct|induced} by inner automorphisms of $\mathcal E_U^{\bar L}$, induced by

\margin{134}
